\section{Tilt stability and numerical transport}
\label{sec:tilt}

We now begin the independent proof of
Theorem~\ref{thm:uniform-inequality}, which applies to arbitrary smooth
projective threefolds along a fixed integral ample direction.  The
inequality will first be proved at one volume, then transported to all
positive volumes.  This section supplies the tilt conventions and
deformation facts needed for that transport.  Objects with zero reduced
charge must be retained: they affect both finite factor filtrations and
duality at real parameters.

\subsection{Characters, slopes, and the first tilt}

Throughout this section, $X$ is a smooth projective complex threefold,
$H$ is an integral ample divisor class, and $h=H^3$.  All intersections
and Chern characters are numerical.  For a real divisor class $T$, set
$\ch^T(E)=e^{-T}\ch(E)$ and use the normalized coordinates
\begin{equation}\label{eq:normalized-characters}
\begin{gathered}
 r(E)=\ch_0(E),\qquad c_T(E)=\ch_1^T(E),\\
 d_T(E)=\frac{H^2c_T(E)}h,\qquad
 w_T(E)=\frac{H\ch_2^T(E)}h,\qquad
 z_T(E)=\frac{\ch_3^T(E)}h.
\end{gathered}
\end{equation}
The subscript $0$ is reserved for the untwisted character.  Fix
$B_0\in N^1(X)_{\QQ}$ and put $B=B_0+bH$ for $b\in\RR$.
We abbreviate the twist $B_0+bH$ by a symbolic subscript $b$;
when evaluating at $b=0$, we write the subscript $B_0$ explicitly.
The exponential formula gives
\begin{align}
 d_b&=d_{B_0}-br,&
 w_b&=w_{B_0}-b d_{B_0}+\tfrac12b^2r,\label{eq:twist-two}\\
 z_b&=z_{B_0}-b w_{B_0}+\tfrac12b^2d_{B_0}-\tfrac16b^3r.
 \label{eq:twist-three}
\end{align}

For a positive-rank coherent sheaf define
$\mu_{H,B}=d_B/r$, and give every nonzero torsion sheaf slope
$+\infty$.  Let $\mu^+$ and $\mu^-$ be the largest and smallest slopes
in the slope Harder--Narasimhan filtration.  The torsion pair and its tilt
are
\begin{equation}\label{eq:first-tilt}
\begin{gathered}
 \mathcal T_{H,B}=\{U:\mu^-_{H,B}(U)>0\},\qquad
 \mathcal F_{H,B}=\{V:\mu^+_{H,B}(V)\le0\},\\
 \mathcal B_{H,B}=\langle\mathcal F_{H,B}[1],\mathcal T_{H,B}\rangle,
\end{gathered}
\end{equation}
where the zero sheaf belongs to both parts.  This is the heart of a
bounded $t$-structure.  For $a>0$ put
\begin{equation}\label{eq:tilt-numerator}
 n_{a,b}=w_b-\tfrac12a^2r,\qquad
 \nu_{a,b}(E)=
 \begin{cases}n_{a,b}(E)/d_b(E),&d_b(E)>0,\\
 +\infty,&d_b(E)=0.
 \end{cases}
\end{equation}
For a general real twist $B$, the same formula defines
$\nu_{a,B}$ using $d_B,w_B$ and $n=w_B-a^2r/2$; the notation
$\nu_{a,b}$ specifies the line $B=B_0+bH$.
An object is tilt-stable, respectively tilt-semistable, if every proper
nonzero subobject $F\subset E$ in $\mathcal B_{H,B}$ satisfies
$\nu(F)<\nu(E/F)$, respectively $\nu(F)\le\nu(E/F)$.
This is the convention of \cite[Definition~2.3]{BMS2016}.
The physical parameter $\omega=\sqrt3aH$ used in
\cite{BMT2014,BMS2016} multiplies our finite slope by $1/(\sqrt3a)$,
so defines the same stable and semistable objects.

Here are two harmless reductions for the inequality.  Write
$q=H^2B_0/h$ and replace $(B_0,b)$ by $(B_0-qH,b+q)$; the twist
$B$ is unchanged, and the new rational $B_0$ satisfies $H^2B_0=0$.
If $H'=sH$ is very ample, use $a'=a/s$ and $b'=b/s$.
The hearts agree and $\nu'_{a',b'}=\nu_{a,b}/s$.  An estimate with
correction $k'(H')^2$ becomes an estimate with correction
$s^2k'H^2$, and a threshold $a'>R'$ becomes $a>sR'$.
We may therefore assume $H$ very ample and $H^2B_0=0$ while proving
the estimates.  Constants used with duality are chosen for both $B_0$
and $-B_0$.

\subsection{Discriminants and finite-slope filtrations}

The Hodge index theorem gives a useful norm on $N^1(X)_{\RR}$.
Define
\begin{equation}\label{eq:divisor-norm}
 \langle C,D\rangle=\frac{HCD}{h},\qquad
 C_\perp=C-\langle H,C\rangle H,\qquad
 \|C\|^2=\langle H,C\rangle^2-\langle C_\perp,C_\perp\rangle.
\end{equation}
The pairing is negative definite on $H^\perp$, so this is a positive
definite norm.  For a positive-rank torsion-free sheaf $G$ put
\begin{equation}\label{eq:sheaf-discriminants}
\begin{gathered}
 \mu(G)=\frac{d_0(G)}{r(G)},\qquad
 \xi_G=\left\langle\frac{c_0(G)}{r(G)},\frac{c_0(G)}{r(G)}\right\rangle
             -\frac{2w_0(G)}{r(G)},\\
 \delta_G^T=\left(\frac{d_T(G)}{r(G)}\right)^2
             -\frac{2w_T(G)}{r(G)}.
\end{gathered}
\end{equation}
Classical Bogomolov--Gieseker gives $\xi_G\ge0$ for slope-semistable
$G$.  Twist invariance of the full discriminant and the orthogonal
decomposition in \eqref{eq:divisor-norm} give
\begin{equation}\label{eq:xi-delta}
 \delta_G^T=\xi_G-\left\langle
 \left(\frac{c_0(G)}{r(G)}-T\right)_\perp,
 \left(\frac{c_0(G)}{r(G)}-T\right)_\perp\right\rangle\ge\xi_G.
\end{equation}
Slope filtrations of torsion-free sheaves may be refined to slope-stable
torsion-free factors: take saturated equal-slope subsheaves inside a
semistable factor and induct on rank.

For arbitrary derived objects the projected discriminant is
\begin{equation}\label{eq:projected-discriminant}
 \Delta(E)=d_b(E)^2-2r(E)w_b(E)
           =d_{B_0}(E)^2-2r(E)w_{B_0}(E).
\end{equation}
Its independence of $b$ follows from \eqref{eq:twist-two}.
We use the following established tilt-stability input for a \emph{fixed}
integral ample class $H$.  There are Harder--Narasimhan filtrations for
all $a>0$ and real $B$; tilt-stability is open in $(a,B)$; the extremal
phases of a fixed derived object vary continuously; and every
tilt-semistable object satisfies
\begin{equation}\label{eq:tilt-discriminants}
 \Delta\ge0,\qquad
 \langle c_B,c_B\rangle-2rw_B+C_Hd_B^2\ge0
\end{equation}
for a constant $C_H\ge0$ depending only on $X,H$.
These are \cite[Theorem~3.5 and Propositions~B.2 and~B.5]{BMS2016},
with volume factors absorbed into $C_H$.  The continuous reduced charge
is $d_B+i(w_B-a^2r/2)$, with heart phase interval
$(-1/2,1/2]$ and phase $1/2$ assigned to its zero-charge objects.
No assertion about varying the direction of $H$ is used here;
\cite[Remark~B.6(b)]{BMS2016} expressly distinguishes that question.

\begin{lemma}\label{lem:tilt-zero-denominator}
An object $E\in\mathcal B_{H,B}$ has $d_B(E)\ge0$.
Write $n=w_B-a^2r/2$.
If $d_B(E)=0$, then $n(E)\ge0$, with equality exactly when
$E$ is a zero-dimensional sheaf, including the zero object.
For a finite-slope object $E$, semistability is equivalent to the
condition that every proper nonzero subobject has positive denominator
and slope at most $\nu(E)$.  Every finite-slope quotient of a
finite-slope semistable object has slope at least $\nu(E)$.
\end{lemma}

\begin{proof}
Write $V=\mathcal H^{-1}(E)$ and $U=\mathcal H^0(E)$ for ordinary
cohomology.  The defining torsion pair gives $d_B(V)\le0$ and
$d_B(U)\ge0$, whence $d_B(E)=d_B(U)-d_B(V)\ge0$.
If it is zero, a nonzero $V$ is torsion-free and slope-semistable of
slope zero, whereas $U$ has dimension at most one.  Bogomolov--Gieseker
and Hodge index give $w_B(V)\le0$; effectiveness of the curve cycle
gives $w_B(U)\ge0$.  Thus
\[
 n(E)=w_B(U)-w_B(V)+\tfrac12a^2r(V)\ge0.
\]
Equality forces $V=0$ and the curve cycle of $U$ to vanish, so $U$ is
zero-dimensional.  The converse is immediate.

In a short exact sequence $0\to F\to E\to G\to0$ with $d(E)>0$,
a subobject with $d(F)=0$ has slope $+\infty$ while $d(G)>0$.
It is forbidden by semistability.  If both denominators are positive,
additivity makes $\nu(E)$ their weighted average, proving the desired
equivalences.  If $d(G)=0$, then $n(G)\ge0$, so
$\nu(F)\le\nu(E)$ and comparison with $\nu(G)=+\infty$ is automatic.
\end{proof}

Strict comparison with the quotient is significant: a zero-dimensional
quotient does not destroy stability in our convention, although it
gives equality between the slopes of the subobject and the object.
An object stable for the stricter subobject--object convention is
semistable in ours, and is therefore covered by all the inequalities
below.

\begin{lemma}\label{lem:tilt-support-bound}
For fixed $a>0$ there is a constant $C_a$, depending only on $a,X,H$,
such that every tilt-semistable object, at any real twist $B$, satisfies
\begin{equation}\label{eq:tilt-support-bound}
 |r|+|w_B|+\|c_B\|\le C_a(d_B+|n|),\qquad n=w_B-a^2r/2.
\end{equation}
The constants can be bounded uniformly when $a$ ranges in a compact
subinterval of $(0,\infty)$.
\end{lemma}

\begin{proof}
The first inequality of \eqref{eq:tilt-discriminants} gives
$a^2r^2+2rn\le d^2$, and hence
$|r|\le 2|n|/a^2+d/a$.  The equality $w=n+a^2r/2$ bounds $|w|$.
Finally the second inequality gives
\[
 \|c_B\|^2=2d^2-\langle c_B,c_B\rangle
 \le(2+C_H)d^2-2rw,
\]
which proves the norm bound.  These expressions also prove uniformity
on compact $a$-intervals.
\end{proof}

This estimate controls rank, the full first character, and one degree
of the second character.  It is not a support statement on the full
numerical Grothendieck group.  Its immediate use is to make a
finite-slope factorization terminate even when $b$ is irrational.

\begin{lemma}\label{lem:finite-tilt-factors}
Every nonzero finite-slope tilt-semistable object has a finite
filtration in $\mathcal B_{H,B}$ whose factors are tilt-stable of the
same finite slope.  Extensions of finite-slope semistable objects of
the same slope are semistable of that slope.
\end{lemma}

\begin{proof}
Let $\nu(E)=u\in\RR$.  If $E$ is not stable, there is a proper
nonzero subobject $F$ with $\nu(F)=\nu(E/F)=u$ and positive
denominators.  The object $F$ is semistable by
Lemma~\ref{lem:tilt-zero-denominator}.  To obtain a splitting into
two semistable objects, we may have to absorb a zero-dimensional
subobject of the quotient back into $F$.  Put $G=E/F$ and take its
Harder--Narasimhan filtration.  Its last factor is a quotient of $E$,
so, if its slope is finite, that slope is at least $u$.  The other
finite slopes are no smaller.  There must be a finite factor since
$d(G)>0$.  Additivity of $n-ud$, together with the nonnegativity of
$n$ on infinite-slope factors, now forces all finite slopes to equal
$u$ and all infinite-slope factors to have $d=n=0$.
Thus $G$ is either semistable or has a zero-dimensional initial
subobject $T$ followed by a semistable quotient of slope $u$.
Replace $F$ by the preimage of $T$.  It is still a proper
equal-slope subobject of $E$, hence semistable, and its quotient is
now semistable as well.

Repeated splitting cannot produce arbitrarily many nonzero factors.
For a semistable factor of slope $u$, Lemma~\ref{lem:tilt-support-bound}
gives $|r|\le C_a(1+|u|)d$.  If $r\ne0$, integrality of rank
therefore bounds $d$ below by $1/(C_a(1+|u|))$.
If $r=0$, then $d=H^2\ch_1/h$ is positive and belongs to
$h^{-1}\ZZ$, so $d\ge1/h$.  Each factor has denominator at least
the smaller of these two positive constants, and their denominators
sum to $d(E)$.  A finite number of splits consequently gives stable
factors.

For the extension assertion, intersect a subobject with the first
term and take its image in the last term.  A nonzero intersection or
image has positive denominator and slope at most $u$ by semistability.
Their sum has the same properties, proving the subobject criterion.
\end{proof}

\subsection{Semistability on the numerical semicircle}

The next observation brings a nonzero finite slope to slope zero.
It also supplies the side points used in the fixed-apex argument.
It extends the calculation of \cite[Lemma~4.3]{BMS2016} to a fixed
transverse twist $B_0$.

\begin{lemma}\label{lem:tilt-semicircle}
Suppose $E$ is tilt-semistable of finite slope $u$ at $(a,b)$.
Put $c=b+u$ and $R=\sqrt{a^2+u^2}$.  Then $E$ is
tilt-semistable at every point
\begin{equation}\label{eq:tilt-semicircle}
 b'=c+x,\qquad a'=\sqrt{R^2-x^2},\qquad -R<x<R,
\end{equation}
and its slope there is $-x$.  In particular it is semistable of
slope zero at $(R,c)$.
\end{lemma}

\begin{proof}
Since $w_b=ud_b+a^2r/2$, the twist formulas give
$w_c=R^2r/2$ and $d_c=d_b-ur$.  Thus
$\Delta=d_c^2-R^2r^2\ge0$.  If $r\ne0$, the negative alternative
$d_c\le-R|r|$ would give
$d_b=d_c+ur\le-(R-|u|)|r|<0$.  It follows that
$d_c\ge R|r|$; for $r=0$ we have $d_c=d_b>0$.
On the whole open arc, therefore,
\[
 d_{b'}=d_c-xr>0,\qquad
 n_{a',b'}=-x(d_c-xr).
\]

Consider the set of points on the arc where $E$ is semistable in
the corresponding heart.  It contains the original point.  It is
closed within the arc: the formal charge has positive real part,
and continuity of the largest and smallest phases makes both
converge to its finite phase.  In particular, this argument also
ensures membership in the limiting heart.
It is open within the arc as well.  At a semistable point take the
finite stable filtration from Lemma~\ref{lem:finite-tilt-factors}.
Every factor has the same slope there and thus the same $c,R$ in
\eqref{eq:tilt-semicircle}.  Openness of stability keeps all these
finitely many factors stable in nearby hearts, with their common
slope $-x$.  Their fixed triangles are short exact sequences in
those hearts, and their extensions are semistable by the preceding
lemma.  Connectedness of the arc proves the assertion.
\end{proof}

\subsection{Duality, including the point defect}

To treat negative rank we use the shifted derived dual
$\mathbb D(E)=R\mathcal Hom(E,\OO_X)[1]$.

\begin{lemma}\label{lem:tilt-duality}
Let $a>0$, let $B$ be real, and let $E\in\mathcal B_{H,B}$
be tilt-semistable of finite slope $u$.  There are a zero-dimensional
sheaf $T_0$ and a distinguished triangle
\begin{equation}\label{eq:tilt-dual-triangle}
 \widetilde E\longrightarrow\mathbb D(E)\longrightarrow T_0[-1]
 \longrightarrow\widetilde E[1],
\end{equation}
where $\widetilde E\in\mathcal B_{H,-B}$ is tilt-semistable of
slope $-u$.  At twists $B$ and $-B$ respectively their coordinates
are related by
\begin{equation}\label{eq:tilt-dual-characters}
 (r,d,w,z)(\widetilde E)
   =\bigl(-r(E),d_B(E),-w_B(E),z_B(E)+\operatorname{length}(T_0)/h\bigr).
\end{equation}
Moreover $\widetilde E=\tau^{\le0}\mathbb D(E)$ for the ordinary
cohomological truncation.
\end{lemma}

\begin{proof}
For rational $\omega=\sqrt3aH$ and rational $B$,
\cite[Proposition~5.1.3(b)]{BMT2014} gives the triangle and
semistability directly: its hypothesis is finite maximal tilt slope,
which holds for a finite-slope semistable object.  In that triangle
$\widetilde E$ has ordinary cohomology in degrees $-1,0$, whereas
$T_0[-1]$ is in degree $1$.  Hence $\widetilde E$ is the stated
ordinary truncation; it depends only on $E$.

First suppose $E$ is stable at real parameters.  By openness it is
stable at a sequence of nearby parameters with rational $\omega,B$.
The rational triangles all have the same truncation
$\tau^{\le0}\mathbb D(E)$ and the same sheaf
$\mathcal H^1(\mathbb D(E))$.  Their mirror finite phases converge
to the phase of slope $-u$, with the denominator still positive.
Continuity of extremal phases therefore gives semistability of this
fixed truncation in $\mathcal B_{H,-B}$.

For a semistable $E$, use Lemma~\ref{lem:finite-tilt-factors} and
induct on its number of stable factors.  For an extension
$0\to E_1\to E\to E_2\to0$, duality reverses the triangle.
We must show that its degree-zero cohomology in the mirror tilt heart
is semistable and that its degree-one cohomology is zero-dimensional.
The cohomology sequence is
\[
 0\longrightarrow\widetilde E_2\longrightarrow
 \mathcal H^0_{\mathcal B_{H,-B}}(\mathbb D(E))
 \longrightarrow\widetilde E_1\longrightarrow T_{0,2}
 \longrightarrow\mathcal H^1_{\mathcal B_{H,-B}}(\mathbb D(E))
 \longrightarrow T_{0,1}\longrightarrow0,
\]
and all other tilt cohomology vanishes.  Subobjects and quotients
of a zero-dimensional sheaf in this heart are zero-dimensional:
additivity first forces both denominators to vanish, then both
nonnegative numerators to vanish, and
Lemma~\ref{lem:tilt-zero-denominator} applies.
Consequently the kernel of $\widetilde E_1\to T_{0,2}$ has the
same positive denominator and the same slope as $\widetilde E_1$;
it is semistable by the subobject criterion.  The degree-zero term
is an extension of it by $\widetilde E_2$, and is semistable of
slope $-u$.  The degree-one term is zero-dimensional.  This proves
\eqref{eq:tilt-dual-triangle}, including its truncation description,
for the extension.

Finally the shifted dual has characters $(-r,c_B,-\ch_2^B,\ch_3^B)$
at twist $-B$.  The term $T_0[-1]$ has third character
$-\operatorname{length}(T_0)$, so taking characters of the triangle
gives \eqref{eq:tilt-dual-characters}.
\end{proof}
